\section{The entity lens}
\label{sec:lens}
\enote{sg}{I think the whole CEL stuff is something we could move to the appendix, so the reader can focus on NEL which now seems to be the main thing.}


\citet{gurnee2026workspace}'s Jacobian lens is broadly
similar to LRE.  But
instead of a relation-specific prompt and a focused target
position (the object),
it computes the average 
$W = \operatorname{mean}_i J_i$ over many diverse prompts and many
output positions. The Jacobian lens also does not use the
``anchor'', the zero-order term of the Taylor expansion; so it
is simply $W = \operatorname{mean}_i J_i$.

We derive our \emph{contrastive entity lens} from
Eq.~\ref{eq:beta}, restated as:
\begin{equation}
  f(s) \approx
\beta W s + b_\beta \;=\;
\operatorname{mean}_i\!\left[\,f(s_i)\right]
+ \beta\operatorname{mean}_i\!\left[\, J_i (s - s_i)\,\right]
\end{equation}
Following \citet{gurnee2026workspace}, we also drop the anchor
term $\operatorname{mean}_i\!\left[\,f(s_i)\right]$; see discussion in 
\S\ref{sec:swap}.
This also means that we do not need $\beta$ since the
scaling does not change the output distribution over
tokens. The final form of the contrastive entity lens is then:
\begin{equation}
\mbox{CEL}(e) = \operatorname{mean}_i\!\left[\, J_i (s_e - s_i)\,\right]
\label{eq:cel}
\end{equation}
CEL takes into account the difference between $s_e$ and
the reference point $s_i$, as it should if we want
to approximate $f(s_e)$ using the first derivative. 
Our original motivation was that we are not interested in what the model associates
with an entity like ``Russia'' in general (e.g.,
``Europe''). Instead we are interested in what the model
associates with ``Russia'', but not with other reference
entities like ``Spain'' and ``Italy''. Examples of such
contrastive associations for ``Russia'' are ``vodka'' and ``Kremlin''.
Our hypothesis was that such top-ranked contrastive associations
give rise to linearity.

In initial experiments, we found that the contrastive entity
lens is quite sensitive to the reference set. For example,
if $s_e$ is a French person and a Spanish person is the
reference entity $s_i$, then the vector difference in 
Eq.~\ref{eq:cel} will enhance ``French''. But if $s_i$ is
another French person, then ``French'' will be demoted.
This is bad for answering questions about $s_e$'s
nationality and native language. We therefore also consider
the following noncontrastive entity lens, which drops the
reference set differences:
\begin{equation}
\mbox{NEL}(e) = \operatorname{mean}_i\!\left[\, J_i s_e \,\right]
\label{eq:nel}
\end{equation}
In what follows,
EL refers to both NEL and CEL.

% and
%``entity lens'' refers to NEL unless
%further specified.

To instantiate EL for factual
relations, we need (i) a prompt that evokes factual associations
and (ii) a reference set of entities for computing the Jacobians $J_i$
(both NEL and CEL) and for serving as the $s_i$ in the
vector differences $s_e - s_i$ (for CEL).

As the prompt we adopt \texttt{"\{\} is known for"}. This is a neutral prompt, i.e.,
it is not customized for any specific factual relation. We
use it for all our  experiments unless otherwise stated.

As the reference set, we use entities of the same type as
the subject, e.g., other countries if the subject is a
country, other landmarks if the subject is a landmark
etc.
NEL applied to an input $s$ outputs anything the model is likely
to say in a context in which an entity of the same type as $s$ occurs, including features
it shares with other same-type entities.
CEL outputs only what differentiates
$s$ from other same-type entities.

\paragraph{Computation.} We do not compute $J_i$ explicitly, but obtain each
Jacobian--vector product $J_i u$ by the double-backward
construction \citep{pearlmutter1994fast}, at the cost of a constant
number of forward and backward passes (Appendix~\ref{sec:jvp}).


%; this is what makes
%Eq.~\ref{eq:cel} affordable at $d_{\text{model}}=4096$

\enote{hs}{is this clearly stated in this section? or do we
  need to repeat it here?
\paragraph{What makes $A_r$ relation-agnostic.} Eq.~\ref{eq:cel} is not specific to $r$: the prompt is neutral and shared,
the layer is fixed at 17 for all relations
(Appendix~\ref{sec:layer17}) and
no (subject,~object) pair of the
relation is used in estimation.
}



